% Independent proof module, 18 November 2026.
% This file is a review artifact, not an edit to an owner manuscript.
% The leave-one-out strategy is motivated by Hogsgaard--Rebeschini--Wegel,
% arXiv:2607.02247 (2026). Their stated theorem is for finite dictionaries
% and bounded losses. The arbitrary-prior and conditional sub-Gaussian
% extension below is derived here; it should receive independent checking.

\section{A sharp expectation oracle for the ordinary final posterior}

Let $\pi$ be any probability prior on jointly measurable functions
$f_w:\mathcal X\to[-1,1]$. Suppose the observations are IID,
$f_0(x)=\mathbb E[Y\mid X=x]\in[-1,1]$, and
$\xi=Y-f_0(X)$ is conditionally sub-Gaussian with proxy $\sigma^2$.
Set $s=1+\sigma^2$ and fix $T\ge256s$.
For a sample $S_n$ of size $n$, let
\[
 dQ_{S_n}(w)=Z_{S_n}^{-1}
   \exp\left\{-\frac1T\sum_{i=1}^n(f_w(X_i)-Y_i)^2\right\}d\pi(w),
 \qquad \bar f_{S_n}=\int f_w\,dQ_{S_n}(w).
\]
This is the ordinary posterior with Gaussian observation variance $T/2$,
using all $n$ observations. No averaging over sample sizes is performed.

\begin{theorem}[Sharp all-data expectation bound]
\label{thm:all-data}
For every deterministic comparison law $\rho\ll\pi$,
\[
 \mathbb E_{S_n} R(\bar f_{S_n})
 \le \int R(f_w)\,d\rho(w)
        +\frac{T\operatorname{KL}(\rho\|\pi)}{n+1}.
\]
Consequently the infimum over all such comparison laws may be taken on
the right-hand side. The theorem requires neither convexity of the
support of $\pi$ nor a Bernstein condition relative to the comparator.
\end{theorem}

\begin{proof}
We first prove an averaged one-observation inequality. Fix a probability
law $q$ on $U\in[-1,1]$, with mean $\mu$ and variance $v$. For a real
label $y$, let $L_y(u)=(u-y)^2$ and
\[
 dp_y(u)=\frac{e^{-L_y(u)/T}}{\int e^{-L_y/T}\,dq}\,dq(u).
\]
For two independent draws $U,U'$ from $q$, the covariance identity gives
\begin{align*}
 \int L_y\,dq-\int L_y\,dp_y
 &=\frac{\mathbb E[(L_y(U)-L_y(U'))
         (e^{-L_y(U')/T}-e^{-L_y(U)/T})]}{2\int e^{-L_y/T}\,dq}.
\end{align*}
Since
\[
 |L_y(u)-L_y(u')|\le2(1+|y|)|u-u'|,
 \qquad \operatorname{osc}_{[-1,1]}L_y\le4(1+|y|),
\]
the mean-value theorem bounds the last display by
\[
 \frac{4(1+|y|)^2}{T}
       e^{4(1+|y|)/T}v.
\]
Also $\int L_y\,dq=L_y(\mu)+v$. Hence
\begin{equation}
 \int L_y\,dp_y\ge L_y(\mu)+
 \left[1-\frac{4(1+|y|)^2}{T}e^{4(1+|y|)/T}\right]v.
 \label{eq:one-step-pointwise}
\end{equation}

Now let $Y=a+\xi$, where $|a|\le1$ and $\xi$ is centered
$\sigma^2$-sub-Gaussian, independently of the fixed law $q$.
The elementary estimates
\[
 \mathbb E e^{t|\xi|}\le2e^{\sigma^2t^2/2},
 \qquad \mathbb E\xi^4\le16\sigma^4
\]
and Cauchy--Schwarz imply
\begin{align*}
 &\mathbb E\left[(1+|Y|)^2e^{4(1+|Y|)/T}\right]\\
 &\quad\le e^{8/T}
       \mathbb E[(8+2\xi^2)e^{4|\xi|/T}]\\
 &\quad\le e^{8/T}
       \left[16e^{8\sigma^2/T^2}
              +8\sqrt2\,\sigma^2e^{16\sigma^2/T^2}\right]
 \le32s
 \qquad(T\ge256s).
\end{align*}
Thus the expected coefficient subtracted in
\eqref{eq:one-step-pointwise} is at most $128s/T\le1/2$.
After averaging over $Y$ we obtain
\begin{equation}
 \mathbb E_Y L_Y(\mu)
 \le\mathbb E_Y\int L_Y\,dp_Y-\tfrac12v
 \le\mathbb E_Y\int L_Y\,dp_Y.
 \label{eq:averaged-one-step}
\end{equation}
The same statement holds conditionally on any input $X=x$ and on an
independent training sample.

Introduce $n+1$ IID observations, retaining the same fixed architecture,
prior and $T$ used by the $n$-observation learner. Let $Q^+$ be their full
posterior and $Q^{-i}$ the posterior after removing observation $i$.
Conditional on the other observations and $X_i$, push $Q^{-i}$ forward
through $w\mapsto f_w(X_i)$ to obtain $q$. It is independent of $Y_i$
given this conditioning. Updating $q$ by the omitted squared-loss factor
gives exactly the corresponding pushforward of $Q^+$.
Applying \eqref{eq:averaged-one-step} and then averaging yields
\[
 \mathbb E (\bar f^{-i}(X_i)-Y_i)^2
 \le\mathbb E\int(f_w(X_i)-Y_i)^2\,dQ^+(w).
\]
Sum over $i=1,\ldots,n+1$. By exchangeability and independence, the
left-hand average is $\mathbb E_{S_n}R(\bar f_{S_n})$.
For each realized enlarged sample the Gibbs variational inequality gives
\begin{align*}
 &\int\sum_{i=1}^{n+1}(f_w(X_i)-Y_i)^2\,dQ^+(w)\\
 &\quad\le\int\sum_{i=1}^{n+1}(f_w(X_i)-Y_i)^2\,d\rho(w)
             +T\operatorname{KL}(\rho\|\pi),
\end{align*}
where the nonnegative term $T\operatorname{KL}(Q^+\|\pi)$ was dropped.
Take expectations and divide by $n+1$ to prove the theorem.
All integrals are finite: the outputs are bounded and the labels have
finite second moment. The posterior normalizers are positive, and the
loss oscillation over parameters is finite for every realized sample.
\end{proof}

\begin{corollary}[Logarithmic comparison entropy]
Suppose a deterministic $g:\mathcal X\to[-1,1]$ has, for every $0<r\le1$,
a data-independent comparison law with
\[
 \mathbb E_\rho\|f_w-g\|_{L^2(P_X)}^2\le r,
 \qquad \operatorname{KL}(\rho\|\pi)\le K_0\log(K_1/r),
 \qquad K_0\ge1,\ K_1\ge e.
\]
At $T=256s$ the same all-data final posterior mean satisfies
\[
 \mathbb E R(\bar f_{S_n})
 \le R(g)+\frac{5+256sK_0\log(K_1(n+1)^2)}{n+1}.
\]
Thus the oracle is sharp, with coefficient one on comparator risk and
a remainder of order $sK_0\log(K_1n)/n$.
\end{corollary}
\begin{proof}
Write $A(g)=\mathbb E(g-f_0)^2\le4$. Cauchy--Schwarz gives
\[
 \int R(f_w)\,d\rho
 \le R(g)+2\sqrt{A(g)r}+r.
\]
Take $r=(n+1)^{-2}$ and use the theorem. The approximation increase is
at most $4/(n+1)+1/(n+1)^2\le5/(n+1)$.
\end{proof}

\paragraph{Scope.}
The result is in expectation over the training sample. It does not
assert a sharp high-probability oracle, and an all-data refit is not
being inferred from a holdout result. Every observation enters the
likelihood of the returned posterior. If the comparison constants are
uniform in learner width, so is this bound. One may take the infimum
over all deterministic comparator classes and resources directly;
unlike a common high-probability event, this expectation oracle needs
no confidence allocation over classes. The learner must still have
its architecture and prior fixed before observing the data.

\paragraph{Attribution.}
The leave-one-out strategy and a pointwise inverse-tilting inequality
for bounded losses are from M.~M.~H\o gsgaard, P.~Rebeschini and
T.~Wegel, \emph{Aggregation with Exponential Weights is Optimal in
Expectation}, arXiv:2607.02247, July 2026. Their main statement uses a
finite uniform dictionary and bounded losses. Above, sums are replaced
by integrals and the variational comparison is written in KL form;
the pointwise bounded-loss step is replaced by an averaged cooling
inequality for bounded predictions and conditionally sub-Gaussian
labels. This extension should not be attributed verbatim to that paper.
